\section{Forecasting and alerting}\label{sec:forecasting}

Community-security organisations and public agencies need to know, on a given day,
whether the coming days are likely to bring an unusual number of incidents. This
section develops four pieces:
\begin{itemize}
  \item predictive distributions from the fitted model;
  \item an alert rule with an explicit trade-off between the costs of false alarms and
  missed surges;
  \item a correction of the model's state for reporting delays;
  \item methods to evaluate forecasts and alerts.
\end{itemize}
Self-exciting models have been used for short-term crime forecasting
\citep{mohler2011,reinhart2018} and for forecasting the popularity of online content
\citep{zhao2015,rizoiu2017}. The decision rule (\cref{thm:alert}(a)) and the scoring
rules are classical. The monotone coupling behind \cref{thm:alert}(b)--(c) is the
standard cluster, or Poisson-embedding, argument \citep{hawkes1974,bremaud1996}. What we
add is its consequence for the shape of the Bayes alert region and the delay-corrected
state with its bias (\cref{prop:nowcast}). Proofs are in \cref{app:attrproofs}; exact
predictive distributions, an imputation algorithm for alerting from delayed data and a
monitor for changes in the parameters are in \cref{app:fc-supp}. Throughout, parameters
are treated as known. Parameter uncertainty enters through posterior or bootstrap draws.

\subsection{Predictive distributions}\label{sec:fc-predictive}

The basic fact is that, given the past, the future is again a Poisson cluster process.
Future salient events are unknown at time $t$. We forecast conditionally on a given
calendar of future events, which is empty unless stated. If event arrivals are modelled
as a marked Poisson process, the results below hold conditionally on its realisation,
and therefore also after averaging over it.
The choice of future calendar matters for calibration. A forecaster that sets the
future calendar to empty while events in fact keep arriving is misspecified: its
predictive law is not the conditional law of the count, and it under-forecasts. A
forecaster that averages over a correct model for event arrivals gives the conditional
law, so its probability integral transform is uniform when windows are pooled
(\cref{lem:fc-pit}). What no forecaster can do is \emph{time} new events, which limits
the recall of an alert rule rather than its pooled calibration (\cref{sec:sim-alert}).

\begin{lemma}[The future given the past]\label{lem:fc-future}
Assume \cref{ass:th-model}, and fix $t\ge0$ and the calendar of events. For $s>t$, let
$\bm u_t(s)=\sum_e\bm\nu_e(s)$ be the remaining immigrant intensity of all events in the
calendar, and let $\bm h_t(s)=\int_{[0,t]}\Phik(s-r)\,\dN(r)$. Conditionally on
$\Hist_t$, the incidents in $(t,\infty)$ form a Poisson cluster process. Its immigrants
in stream $k$ form a Poisson process with intensity
\[
\nu_{t,k}(s)=\mu_k(s)+u_{t,k}(s)+h_{t,k}(s),\qquad s>t,
\]
and offspring are generated as in \cref{sec:th-cluster}. If $\spr(\Amat)<1$, then
$\E[\bm N(t,t+\Delta_{\mathrm a}]\mid\Hist_t]=\int_t^{t+\Delta_{\mathrm a}}(\mathcal K_t\bm\nu_t)(s)\,\dd s$,
with $\mathcal K_t$ as in \eqref{eq:at-cascade}.

Suppose in addition that $\varphi_{kj}(u)=\alpha_{kj}\beta_{kj}e^{-\beta_{kj}u}$ and
$\psi_k(u)=\eta_ke^{-\eta_ku}$, and that there are no future events. Then, for $s>t$,
\begin{equation}\label{eq:fc-state}
\begin{aligned}
  h_{t,k}(s)&=\sum_{j=1}^{M}x_{kj}(t)\,e^{-\beta_{kj}(s-t)}, &
  x_{kj}(t)&=\sum_{i:\,t^j_i\le t}\alpha_{kj}\beta_{kj}e^{-\beta_{kj}(t-t^j_i)},\\
  u_{t,k}(s)&=u_k(t)\,e^{-\eta_k(s-t)}, &
  u_k(t)&=\sum_{e:\,\tau_e\le t}\upsilon_kS_e\eta_ke^{-\eta_k(t-\tau_e)} .
\end{aligned}
\end{equation}
The state $\bm\chi(t)=(\bm x(t),\bm u(t))\in[0,\infty)^{M^2+M}$ is a piecewise-deterministic
Markov process. Between incidents and events it decays exponentially. At a stream-$j$
incident, $x_{kj}$ jumps by $\alpha_{kj}\beta_{kj}$ for each $k$. At event $e$, $u_k$
jumps by $\upsilon_kS_e\eta_k$. The conditional law of the future given $\Hist_t$ depends on
$\Hist_t$ only through $\bm\chi(t)$, and the predictive mean is affine in $\bm\chi(t)$.
The intensity just after $t$ is
$\lambda_k(t+)=\mu_k(t)+u_k(t)+\sum_jx_{kj}(t)$.
\end{lemma}

\Cref{lem:fc-future} makes exact forecasting possible, with no discretisation error.
Generation~0 of the future consists of background immigrants on
$(t,t+\Delta_{\mathrm a}]$, the remaining immigrants of past events, and the pending
offspring of past incidents; those of stream $j$ in stream $k$ number
$\Pois\{x_{kj}(t)(1-e^{-\beta_{kj}\Delta_{\mathrm a}})/\beta_{kj}\}$. Each generation-0 point
then cascades as in \cref{sec:th-cluster}, truncated at $t+\Delta_{\mathrm a}$, and $J$
repetitions give a Monte Carlo predictive distribution with error of order $J^{-1/2}$.
The predictive law can also be computed without simulation. Its probability generating
function solves $M^2$ ordinary differential equations backwards in time, and a discrete
Fourier transform then gives the predictive probability mass function
(\cref{rem:fc-pgf} in \cref{app:fc-pgf}).

\subsection{Bayes-optimal alerting}\label{sec:fc-alert}

Fix a window length $\Delta_{\mathrm a}$ and a \emph{surge set} $\mathcal U\subseteq\mathbb Z_+^M$
that is an upper set: if $\bm n\in\mathcal U$ and $\bm n'\ge\bm n$, then
$\bm n'\in\mathcal U$. The leading example is $\mathcal U=\{\bm n:n_k\ge c\}$, at least
$c$ incidents in stream $k$. Weighted totals $\{\bm n:\bm w^{\top}\bm n\ge c\}$ with
$\bm w\ge\zero$ also qualify. The surge indicator is
$Y_t=\ind{\bm N(t,t+\Delta_{\mathrm a}]\in\mathcal U}$.

At time $t$ the analyst holds information $\mathcal G_t$. This is the history $\Hist_t$
when data are complete and timely, and less under reporting delays
(\cref{sec:fc-nowcast}). The analyst chooses an alert $a_t\in\{0,1\}$ that is
$\mathcal G_t$-measurable, and incurs the loss
$c_{\mathrm{FA}}\,a_t(1-Y_t)+c_{\mathrm{miss}}(1-a_t)Y_t$ with
$c_{\mathrm{FA}},c_{\mathrm{miss}}>0$. Let
\[
\pi_t=\Pr(Y_t=1\mid\mathcal G_t),\qquad
\pi^*=\frac{c_{\mathrm{FA}}}{c_{\mathrm{FA}}+c_{\mathrm{miss}}} .
\]

\begin{theorem}[Bayes-optimal alerting]\label{thm:alert}
\begin{enumerate}[(a)]
  \item \emph{Bayes rule.} Among all $\mathcal G_t$-measurable alerts, the rule
  $a^*_t=\ind{\pi_t\ge\pi^*}$ minimises the expected loss. Its risk is
  $\E\min\{c_{\mathrm{FA}}(1-\pi_t),\,c_{\mathrm{miss}}\pi_t\}$. Every other minimiser
  agrees with $a^*_t$ almost surely outside $\{\pi_t=\pi^*\}$.
  \item \emph{Monotonicity in the history.} Let $\mathcal G_t=\Hist_t$, and fix the
  parameters and the calendar of events. Let $\omega\subseteq\omega'$ be two histories
  on $[0,t]$, where $\omega'$ contains every incident of $\omega$, with the same stream
  labels, and possibly more. Then there is a coupling of the predictive laws of
  $\bm N(t,t+\Delta_{\mathrm a}]$ under $\omega$ and $\omega'$, $\bm N$ and $\bm N'$, with
  $\bm N'\ge\bm N$. In particular $\pi_t(\omega)\le\pi_t(\omega')$ for every upper
  set $\mathcal U$.
  \item \emph{Exponential kernels: a monotone decision region.} Under the
  exponential specification of \cref{lem:fc-future}, $\pi_t=G_t(\bm\chi(t))$ for a
  function $G_t$ that is nondecreasing in each coordinate and Lipschitz. The Bayes
  alert region $\mathcal A_t=\{\bm\chi:G_t(\bm\chi)\ge\pi^*\}$ is therefore a closed
  upper set: if the rule alerts at $\bm\chi$, it alerts at every $\bm\chi'\ge\bm\chi$.
  The function $G_t$ depends on $t$ only through $\bm\mu$ on $(t,t+\Delta_{\mathrm a}]$ and the
  calendar of future events. For constant $\bm\mu$ and no future events it does not
  depend on $t$.
  \item \emph{Univariate case: a threshold curve.} Let $M=1$, with $\mu>0$ on
  $(t,t+\Delta_{\mathrm a}]$ and $\mathcal U=\{n\ge c\}$ for some $c\ge1$. Then
  $\pi_t=g_t(x(t),u(t))$, where $g_t$ is continuous, nondecreasing in $u$ and strictly
  increasing in $x$. The Bayes rule alerts if and only if
  \[
  x(t)\ \ge\ x^*_t(u(t)),\qquad x^*_t(u)=\inf\{x\ge0:g_t(x,u)\ge\pi^*\}<\infty .
  \]
  The threshold curve $x^*_t(u)$ is nonincreasing in $u$, nondecreasing in $c$ and in
  $\pi^*$, and nonincreasing in $\Delta_{\mathrm a}$. If $\eta=\beta$, or if no event has occurred
  so that $u\equiv0$, then $g_t$ depends on $x+u$ only. The rule is then a threshold on
  the intensity: alert if and only if $\lambda(t+)\ge\lambda^*_t:=\mu(t)+x^*_t(0)$.
  In general, replacing $u(t)$ by $0$ changes $\pi_t$ by at most $L_tu(t)$, where $L_t$
  is a Lipschitz constant of $G_t$ in (c), for instance $L_t=\Delta_{\mathrm a}/(1-n)$.
\end{enumerate}
\end{theorem}

Part (a) is standard Bayes decision theory. Part (b) says that the predictive law of
$\bm N(t,t+\Delta_{\mathrm a}]$ increases with the history in the usual multivariate stochastic
order \citep{shaked2007}. The decision depends on the two costs only
through $\pi^*$, which is agreed with partners. The analysis plan's grid of cost ratios
$c_{\mathrm{miss}}/c_{\mathrm{FA}}\in\{2,5,10\}$ corresponds to
$\pi^*\in\{1/3,1/6,1/11\}$ (\cref{sec:plan-estimands}).

The structural parts have three operational consequences. Under the model, an
additional incident never turns an alert off. In the univariate case, the curve
$x^*_t(u)$ can be tabulated once, and checking for an alert means comparing the current
excitation with it; since an event's contribution $u(t)$ never vanishes, the single
intensity threshold $\lambda^*_t$ is exact only when $\eta=\beta$ or before any event. In
the multivariate case, $\mathcal A_t$ is an upper set but in general not a half-space. The
predictive mean is affine in $\bm\chi(t)$, so the rule ``alert if the predictive mean
exceeds a level'' has a half-space as its region. That rule is not optimal in general,
because the tail of the future count also depends on its dispersion, which is larger when
the excitation sits in a strongly self-exciting stream.
With unknown parameters, part (a) applies to the posterior predictive probability.
Parts (b)--(d) hold for each parameter value, but they need not hold for the posterior
mixture, because the posterior itself depends on the history.

\subsection{Nowcasting under reporting delays}\label{sec:fc-nowcast}

Incidents reach the analyst late (\cref{sec:data-delays}). Near the forecast origin the
record is therefore incomplete. The model's state is then too low, exactly when an alert
matters most.

\begin{assumption}[Recording and reporting]\label{ass:fc-report}
Each stream-$k$ incident is recorded with probability $p_k\in(0,1]$. A recorded incident
that occurs at $s$ becomes available to the analyst at $s+d$, where the delay $d$ has distribution
function $F_k$ on $[0,\infty)$. Recording indicators and delays are independent across
incidents, and independent of the incident process and the calendar of events.
\end{assumption}

Fix the analysis time $t$, and let $O^{(t)}_k$ count the occurrence times of the
stream-$k$ incidents available by $t$. Let $\mathcal G^{\circ}_s$ be generated by
$\Hist_s$ together with the recording indicators and delays of the incidents in
$[0,s]$. We call an incident \emph{available} if it is recorded and has been reported by
$t$.

\begin{proposition}[Nowcasting the state]\label{prop:nowcast}
Assume \cref{ass:th-model,ass:fc-report}.
\begin{enumerate}[(a)]
  \item \emph{Thinning.} On $[0,t]$, $O^{(t)}_k$ has $\mathcal G^\circ$-intensity
  $p_kF_k(t-s)\,\lambda_k(s)$, and the incidents not yet available have
  $\mathcal G^\circ$-intensity $\{1-p_kF_k(t-s)\}\lambda_k(s)$. Hence
  $\E\,O^{(t)}_k(\dd s)=p_kF_k(t-s)\,\E\lambda_k(s)\,\dd s$.
  \item \emph{Unbiased nowcast.} Let $\delta\ge0$ satisfy $F_j(\delta)>0$ for every
  $j$. For $s\ge t-\delta$, define the weighted carried-over excitation
  \begin{equation}\label{eq:fc-ht}
    \hat{\bm h}(s)=\sum_{j=1}^{M}\ \sum_{\substack{i:\ s^j_i\le t-\delta,\\ \text{available}}}
    \frac{\Phik_{\cdot j}(s-s^j_i)}{p_jF_j(t-s^j_i)},
  \end{equation}
  where $\Phik_{\cdot j}$ is the $j$th column of $\Phik$. Then
  $\E[\hat{\bm h}(s)\mid\Hist_{t-\delta}]=\bm h_{t-\delta}(s)$ for every
  $s\ge t-\delta$. If $\spr(\Amat)<1$, the nowcast
  $\hat{\bm\lambda}(t)=\mathcal K_{t-\delta}(\bm\mu+\bm u_{t-\delta}+\hat{\bm h})(t)$
  satisfies
  $\E[\hat{\bm\lambda}(t)\mid\Hist_{t-\delta}]=\E[\bm\lambda(t)\mid\Hist_{t-\delta}]$,
  and hence $\E\hat{\bm\lambda}(t)=\E\bm\lambda(t)$.

  For exponential kernels the same holds for the state. Let $\hat{\bm x}(t)$ be
  obtained by computing the state at $t-\delta$ with the weights of \eqref{eq:fc-ht}
  and propagating it to $t$ with the mean dynamics of \cref{prop:irf}(c). Then
  $\hat{\bm x}(t)$ is conditionally unbiased for $\E[\bm x(t)\mid\Hist_{t-\delta}]$.
  If every $F_j(0)>0$, one can take $\delta=0$. In that case $\hat{\bm h}$ is the
  Horvitz--Thompson estimator of $\bm h_t$, and it is conditionally unbiased given the
  complete history $\Hist_t$.
  \item \emph{The naive state is biased downward.} Let
  $\tilde x_{kj}(t)=\sum_{i:\,s^j_i\le t,\ \text{available}}\varphi_{kj}(t-s^j_i)/p_j$ be
  the state computed from the available incidents, with recording corrected but delays
  ignored. Then
  \begin{equation}\label{eq:fc-naive}
    \E\bigl[\tilde x_{kj}(t)\mid\Hist_t\bigr]
    =\sum_{i:\,s^j_i\le t}F_j(t-s^j_i)\,\varphi_{kj}(t-s^j_i)
    \ \le\ x_{kj}(t)=\sum_{i:\,s^j_i\le t}\varphi_{kj}(t-s^j_i).
  \end{equation}
  The inequality is strict whenever some incident with $\varphi_{kj}(t-s^j_i)>0$ has
  $F_j(t-s^j_i)<1$. Now suppose $\bm\mu$ is constant, there are no events,
  $\spr(\Amat)<1$ and the process is stationary. Then
  $\E\tilde x_{kj}(t)=\Pr(d_j\le U_{kj})\,\E x_{kj}(t)$, where $d_j\sim F_j$ and
  $U_{kj}\sim g_{kj}$ are independent. For exponential kernels,
  $\Pr(d_j\le U_{kj})=\E e^{-\beta_{kj}d_j}$.
  \item \emph{Consequence for alerts.} If $p_k\equiv1$, then
  $\tilde{\bm x}(t)\le\bm x(t)$ pathwise. By \cref{thm:alert}(c), the alert probability
  computed from the naive state is then at most the one computed from the true state. A
  rule based on the naive state alerts only when the rule based on the true state
  alerts. Ignoring delays can postpone an alert but never brings one forward.
\end{enumerate}
\end{proposition}

Part (a) restates \eqref{eq:data-righttrunc} at the level of intensities: any likelihood
for the data available at $t$ must carry the factor $p_kF_k(t-s)$ on $\lambda_k(s)$, and
integrate over the incidents not yet available. Data augmentation, by Markov chain Monte
Carlo or stochastic EM, imputes them, the self-exciting analogue of Bayesian nowcasting
\citep{hohle2014,mcgough2020}. A simpler route, adequate for estimating $\vartheta$,
restricts the likelihood to the settled part of the record, $s\le t-L$ with
$F_k(L)\approx1$, and uses (b) only for the state near $t$.

The size of the bias in (c) can be large. With exponentially distributed delays of mean
two days and a kernel rate $\beta=1.5$ per day, $\E e^{-\beta d}=1/(1+2\beta)=0.25$. On
average, the naive state then carries a quarter of the true excitation.

The unbiased nowcast has a cost in variance. The weights $1/\{p_jF_j(t-s)\}$ are large
for recent incidents, and without a lag floor they are unbounded: if $F_j(0)=0$,
$F_j(u)=O(u)$ as $u\downarrow0$ and $\varphi(0)>0$, the conditional variance of the weighted
sum over all incidents up to $t$ has infinite expectation. This is why (b) requires
$F_j(\delta)>0$; a floor $\delta>0$ discards the available incidents in $(t-\delta,t]$.

An unbiased state does not give a calibrated tail probability. $G_t$ is nonlinear, and
a noisy state that happens to be high produces a high plug-in probability more often
than a surge. Plugging $\hat{\bm x}(t)$ into $G_t$, or simulating forward from it,
therefore does not give $\Pr(Y_t=1\mid\mathcal G_t)$ for the available data
$\mathcal G_t$. The Bayes optimality of \cref{thm:alert}(a) then holds only for the
oracle rule that knows the true state. Plugging a nowcast into the threshold rule is a
shortcut, not the Bayes rule for delayed data. \Cref{alg:alert} in \cref{app:fc-sis}
approximates the predictive probability given the available data instead. It imputes the
incidents that are not yet available, by sequential importance sampling with the
intensities of (a) (\cref{lem:fc-is}), and so targets the Bayes rule under delays. We have
not run it: the simulation study evaluates only the plug-in shortcut
(\cref{sec:sim-alert}). Whichever rule is used, the probability integral transform
(\cref{lem:fc-pit}) and a reliability diagram of the alert probabilities should be checked
on held-out data before operational use. If they show miscalibration, the probabilities
should be recalibrated, for instance by isotonic regression of $Y_t$ on $\hat\pi_t$ over
the training period, before the threshold $\pi^*$ is applied.

\subsection{Evaluation}\label{sec:fc-eval}

Forecasts are evaluated against \emph{settled} counts. These are the counts in each
window once reporting has essentially completed, not the counts available at the
forecast origin. Let $P_t$ be the predictive distribution of a count
$Y_t=N_k(t,t+\Delta_{\mathrm a}]$, with distribution function $P_t(y)=P_t(Y\le y)$.

\paragraph{Proper scoring rules} A scoring rule is strictly proper if the true
predictive distribution uniquely optimises its expected value \citep{gneiting2007}. We
use two scores, both negatively oriented:
\begin{equation}\label{eq:fc-scores}
  \mathrm{LogS}(P,y)=-\log P(Y=y),\qquad
  \mathrm{CRPS}(P,y)=\sum_{z=0}^{\infty}\bigl\{P(z)-\ind{y\le z}\bigr\}^2
  =\E_P|X-y|-\tfrac12\E_P|X-X'| ,
\end{equation}
where $X$ and $X'$ are independent with law $P$. The log score is local and requires
$P(y)>0$, so Monte Carlo probability mass functions need a floor, or the exact law of
\cref{rem:fc-pgf} should be used. The CRPS for counts is strictly proper relative to
distributions with a finite mean, is insensitive to small probabilities, and is estimated
from predictive draws by the sample version of the second expression
\citep{gneiting2007,czado2009}. For the surge event itself, the Brier score
$(\pi_t-Y_t)^2$ is strictly proper.

\paragraph{Calibration} For discrete counts the probability integral transform (PIT)
is randomised:
\[
U_t=P_t(Y_t-1)+V_t\bigl\{P_t(Y_t)-P_t(Y_t-1)\bigr\},
\]
with $V_t$ uniform on $(0,1)$ and independent of everything else, and $P_t(-1)=0$.
\Citet{czado2009} also give a non-randomised version.

\begin{lemma}[PIT for counts]\label{lem:fc-pit}
\begin{enumerate}[(a)]
  \item If $P_t$ is the conditional law of $Y_t$ given the forecaster's information
  $\mathcal G_t$, then $U_t$ is uniform on $(0,1)$ and independent of $\mathcal G_t$.
  \item Suppose forecasts are issued at $t_m=t_0+m\Delta_{\mathrm a}$, so that the windows do not
  overlap. Suppose also that each past count $Y_{t_l}$, $l<m$, is
  $\mathcal G_{t_m}$-measurable, that is, past windows have settled by the next forecast
  origin. Then $U_{t_1},U_{t_2},\dots$ are independent and uniform on $(0,1)$.
\end{enumerate}
\end{lemma}

This is the count-data form of the prequential principle of \citet{dawid1984}; see also
\citet{gneiting2007b} and \citet{czado2009}. A slope in the PIT histogram indicates bias,
a U shape forecasts that are too narrow, and a hump forecasts that are too wide. With
overlapping windows, or with counts that have not settled, the $U_t$ are uniform but
dependent, so tests of uniformity need block resampling.

\paragraph{Alerts} For a threshold $\pi$, the hit rate (recall) is
$\mathrm{HR}(\pi)=\Pr(\pi_t\ge\pi\mid Y_t=1)$, the false-alarm rate
$\mathrm{FAR}(\pi)=\Pr(\pi_t\ge\pi\mid Y_t=0)$ and the precision
$\mathrm{PPV}(\pi)=\Pr(Y_t=1\mid\pi_t\ge\pi)$. The ROC curve, $(\mathrm{FAR},\mathrm{HR})$, does not depend on how common surges are.
When surges are rare, the precision--recall curve, $(\mathrm{HR},\mathrm{PPV})$, is more
informative about the burden of false alarms. The realised cost over a test period,
$\sum_t\{c_{\mathrm{FA}}a_t(1-Y_t)+c_{\mathrm{miss}}(1-a_t)Y_t\}$, is the criterion of
\cref{thm:alert}. The comparison with baselines in the analysis plan
(\cref{sec:plan-decisions}) is made in these units. Alerts issued daily for overlapping
windows have dependent outcomes, so intervals use a block bootstrap.

The Brier score is twice the average realised cost when the normalised cost
$\pi^*$ is drawn uniformly from $(0,1)$ (\cref{rem:fc-brier} in \cref{app:fc-brier}). It
summarises performance when partners cannot agree on one cost ratio; when they can, the
realised cost at $\pi^*$ is the relevant measure.

\paragraph{Regime shifts} Alerts concern the next $\Delta_{\mathrm a}$ days with the
parameters held fixed. Whether the parameters themselves have changed, through a sustained
rise either in the background or in contagion, is a different question.
\Cref{sec:fc-cusum} sketches a CUSUM monitor for it (\cref{rem:fc-cusum}). It inherits the
identification limits of \cref{sec:identification} and must be run on settled data or with
the nowcast of \cref{prop:nowcast}.
